\documentclass[11pt]{article}
\usepackage[margin=1in]{geometry}
\usepackage{amsmath,amssymb,amsthm}
\usepackage[hidelinks]{hyperref}

\title{Literature Status for arXiv:1309.3501 Problems 1 and 2}
\author{FA/Banach automatic math run}
\date{2026-07-03}

\begin{document}
\maketitle

\section*{Status}

\textbf{Status:} literature already answered, for the intended graph
formulation of the intrinsic-metric and algebra-characterization questions.

\section*{Source Question}

The source paper is Agelos Georgakopoulos, Sebastian Haeseler, Matthias
Keller, Daniel Lenz, and Radoslaw K. Wojciechowski, \emph{Graphs of finite
measure}, arXiv:1309.3501 \cite{GHKLW}. In Section 8.3, page 53, the paper
asks:
\begin{quote}
Problem 1. Does $(D)$ imply $(C)$ in general?
\end{quote}
Here $(D)$ is the condition that all finite-energy functions are bounded, also
called canonical compactifiability, and $(C)$ is total boundedness with respect
to every intrinsic metric associated with a finite measure. The same page also
asks:
\begin{quote}
Problem 2. Does the fact that $\widetilde D$ is an algebra imply that $(b,c)$
is canonically compactifiable?
\end{quote}

\section*{Supporting Answer}

The supporting paper is Simon Puchert, \emph{Characterizations of canonically
compactifiable graphs via intrinsic metrics and algebraic properties},
arXiv:1803.01942 \cite{Puchert}. Its abstract states that it answers questions
from \cite{GHKLW}.

For Problem 1, Puchert's Theorem 1 proves that if a graph $(X,b)$ is
canonically compactifiable and $\sigma$ is intrinsic with respect to a finite
measure, then $(X,\sigma)$ is totally bounded. Together with the converse from
\cite{GHKLW}, Corollary 1 states that a graph is canonically compactifiable if
and only if it is totally bounded with respect to every finite-measure
intrinsic metric.

For Problem 2, Puchert's Theorem 2 proves that if the finite-energy space
$D(G)$ is an algebra, then the graph $G$ is canonically compactifiable.
Together with the converse already known from \cite{GHKLW}, Corollary 2 states
that $G$ is canonically compactifiable if and only if $D(G)$ is an algebra.

\section*{Scope}

Puchert formulates the result for graphs $G=(X,b)$, the standard zero-killing
setting of the metric characterization. This is the formulation Puchert
explicitly identifies as answering the questions from \cite{GHKLW}. This note
does not claim that every possible arbitrary-killing-term reading of the
source notation $(b,c)$ is covered.

\section*{Search Bounds}

The target was checked against the run registry, solution index, attempt
index, and proof-gap index for arXiv:1309.3501 and core phrases including
``canonically compactifiable'', ``finite measure intrinsic metric'', and
``finite energy algebra''. Web search for the exact source title and these
phrases found arXiv:1803.01942 as the explicit later answer. No competing
later source was needed for the identification.

\section*{Human Review Recommendation}

Mark Problems 1 and 2 from arXiv:1309.3501 as already answered in later
literature for the intended graph formulation. Future work should only reopen
this target if it explicitly studies a variant involving arbitrary killing
terms.

\begin{thebibliography}{9}

\bibitem{GHKLW}
A. Georgakopoulos, S. Haeseler, M. Keller, D. Lenz, and R. K. Wojciechowski,
\emph{Graphs of finite measure}, arXiv:1309.3501.

\bibitem{Puchert}
S. Puchert, \emph{Characterizations of canonically compactifiable graphs via
intrinsic metrics and algebraic properties}, arXiv:1803.01942.

\end{thebibliography}

\end{document}
